\documentclass[../main]{subfiles}
\begin{document}
\chapter{Aplicación de Calidad al Proyecto}
\actividad{ Actividad para el proyecto}
\begin{enumerate}
  \item  Realizar un diagrama de Ishikawa para el proyecto respetando
    las etapas enunciadas.

    \begin{itemize}
      \item Identificar posibles causas de problemas de calidad en el producto.
      \item Utilizar el Diagrama de Ishikawa para analizar y
        categorizar estas causas.
      \item Desarrollar soluciones para abordar las causas y mejorar la
        calidad del producto.
    \end{itemize}

    \subsection*{Optimización de Procesos de Fabricación:}
    \begin{itemize}
      \item Seleccionar un proceso específico en la fabricación del producto.
      \item Usar el Diagrama de Ishikawa para identificar causas de
        ineficiencia en el proceso.
      \item Proponer soluciones para optimizar el proceso y aumentar la
        eficiencia.
    \end{itemize}

    \subsection*{Mejora de la Seguridad en la Producción:}
    \begin{itemize}
      \item Enfocarse en aspectos de seguridad en el área de producción.
      \item Utilizar el Diagrama de Ishikawa para identificar
        posibles causas de riesgos laborales.
      \item Desarrollar estrategias para mejorar la seguridad y
        prevenir accidentes.
    \end{itemize}

    \subsection*{Desarrollo de Nuevos Productos:}
    \begin{itemize}
      \item Trabajar en grupos para diseñar nuevos productos.
      \item Utilizar el Diagrama de Ishikawa para explorar desafíos
        en el desarrollo del producto.
      \item Identificar causas potenciales de problemas en aspectos
        como funcionalidad y estética.
    \end{itemize}

\end{enumerate}

\end{document}